\chapter{Software, Datasets, and Benchmarks}
\label{ch:software}

\section{Software ecosystems}
\subsection{Julia}
\begin{itemize}
\item \textbf{SciML.jl} (DifferentialEquations.jl, Lux.jl,
DiffEqFlux.jl, ModelingToolkit.jl, NeuralPDE.jl, NeuralOperators.jl):
the premier ecosystem for differentiable scientific computing,
universal differential equations, and symbolic-numeric
modelling~\cite{rackauckas2017diffeq}.
\item \textbf{ModelingToolkit + ModelingToolkitNeuralNets}: symbolic
PDE/DAE modelling with embedded neural blocks.
\item \textbf{Optimization.jl, Lux.jl, Zygote.jl, Enzyme.jl}: forward,
reverse, and source-to-source AD.
\end{itemize}

\subsection{Python / JAX}
\begin{itemize}
\item \textbf{JAX} + \textbf{Diffrax}: differentiable solvers (Tsit5,
KenCarp, etc.) integrated with neural networks via
Equinox/Flax/Haiku.
\item \textbf{JAX-CFD, JAX-Fluids, JAX-MD, JAX-FEM}: fluid, molecular,
and FEM solvers.
\item \textbf{PhiFlow, Modulus}: NVIDIA's PINN/operator framework with
multi-GPU support.
\item \textbf{Neuraloperator (PyTorch)}: official FNO/TFNO
implementation.
\item \textbf{DeepXDE}: PINNs and DeepONets with multi-backend support.
\end{itemize}

\subsection{Other}
\begin{itemize}
\item \textbf{Firedrake / FEniCS / dolfin-adjoint}: differentiable FEM.
\item \textbf{PyTorch Geometric, DGL}: GNN backbones.
\item \textbf{e3nn, ESCNN}: equivariant network libraries.
\item \textbf{TorchDiffEq, TorchSDE}: neural ODE/SDE in PyTorch.
\end{itemize}

\section{Datasets and benchmarks}
\subsection{PDE benchmarks}
\begin{itemize}
\item \textbf{PDEBench}~\cite{takamoto2022pdebench}: advection, Burgers,
Navier--Stokes (compressible/incompressible), Darcy, diffusion-reaction,
shallow water; standardised metrics; widely cited.
\item \textbf{The Well}~\cite{ohana2024well}: 15 TB / 16 datasets across
fluids, MHD, biology, granular media; built explicitly for foundation
operators.
\item \textbf{PDEArena}: clean PyTorch implementations of FNO, U-Net,
DeepONet on Navier--Stokes / shallow water.
\item \textbf{BubbleML}: multiphase flows.
\item \textbf{CFDBench, AirfRANS}: aerodynamics.
\end{itemize}

\subsection{Molecular / materials}
\begin{itemize}
\item \textbf{MD17, rMD17, MD22}: gas-phase MD trajectories.
\item \textbf{OC20, OC22}: catalysis with millions of DFT calculations.
\item \textbf{QM9, QMugs}: quantum-chemistry property prediction.
\item \textbf{Matbench, JARVIS}: materials property prediction.
\end{itemize}

\subsection{Climate / weather}
\begin{itemize}
\item \textbf{ERA5}: ECMWF reanalysis, the de-facto training corpus for
ML weather models.
\item \textbf{WeatherBench 2}~\cite{rasp2024wb2}: standardised
evaluation; benchmarks GraphCast, Pangu, FourCastNet, GenCast,
NeuralGCM, etc.
\item \textbf{ClimateBench, ClimSim}: climate emulation, subgrid
parameterisation.
\end{itemize}

\subsection{Astrophysics}
\begin{itemize}
\item \textbf{CAMELS}: thousands of cosmological simulations with
varying physics for ML training.
\item \textbf{IllustrisTNG, EAGLE}: galaxy-formation simulation suites.
\item \textbf{LSST, Euclid mocks}: synthetic survey data.
\end{itemize}

\section{Reproducibility checklist}
\begin{recipe}{Publishing trustworthy SciML}
\begin{enumerate}
\item Code: open-source repository, fixed seeds, environment files.
\item Data: list source, version, preprocessing, splits.
\item Compute: report wall time, hardware, energy where possible.
\item Baselines: at least one classical solver and one prior NN method.
\item Metrics: error norm, calibration, OOD generalisation.
\item Ablations: which inductive biases matter.
\item Limitations: failure cases, regime of validity.
\end{enumerate}
\end{recipe}

\section{Hardware considerations}
\begin{itemize}
\item GPUs (H100/H200/B200) for transformer/operator training; tensor
cores reward large batched FFTs.
\item TPUs (v5p) for very large foundation models (NeuralGCM,
GraphCast).
\item AMX / Apple Silicon for medium-scale prototyping (Mac mini M4
remains highly capable for FNO and PINN prototyping with MPS or
PyTorch--Metal).
\item Mixed precision (BF16/FP16) for inference; \textbf{always train
PINNs in FP64} (\cref{ch:pinns}).
\end{itemize}
